\section{Background with Reference to Sources}

\subsection{Previous Work in Route-Aware Traffic Prediction}

This research builds upon previous work in traffic prediction and route-aware modeling. Voloch and Voloch-Bloch~\cite{voloch2021} demonstrated real-time future traffic estimations for navigation route optimization, establishing the foundation for dynamic traffic prediction. Their approach used simulation-based methods to estimate future traffic conditions and optimize navigation routes accordingly. Building on this, Voloch et al.~\cite{SmartSimulativeRoute2025} introduced smart simulative route predictions that anticipate future congestion through heuristic simulation, improving ETA accuracy. These works laid the methodological groundwork for graph-based analysis of traffic dynamics and route-aware modeling, demonstrating the value of incorporating route information and future traffic predictions in navigation systems.

\subsection{State-of-the-Art ETA Prediction Methods}

With the growth of urban mobility datasets, machine learning models have been applied to ETA prediction. Tree-based methods such as XGBoost~\cite{chen2016xgboost} leverage handcrafted features (trip distance, time-of-day) and perform well on aggregated trip records, including NYC Taxi~\cite{nyc_tlc} and Porto Taxi~\cite{moreira2013porto} datasets. However, they cannot capture fine-grained spatio-temporal interactions.

Trajectory- and trip-based neural models address this gap. DeepTTE~\cite{deepTTE2018} learns ETA from raw GPS traces and reports results on city-scale taxi/ride-hailing datasets with per-trip horizons typically within tens of minutes. STAD~\cite{abbar2020stad} adjusts routing-engine outputs with spatio-temporal corrections, improving travel-time estimation. Production systems such as STANN~\cite{stann2021} evaluate on large ride-hailing ETA datasets with short-to-medium trip durations.

In parallel, graph-based spatio-temporal learning has advanced traffic prediction by representing road networks as graphs. DCRNN~\cite{dcrnn2018} combined diffusion graph convolutions with recurrent units and became a common backbone for traffic speed/flow prediction on sensor benchmarks (METR-LA~\cite{li2018_metrla}, PEMS-BAY~\cite{pemsbay}). ST-GCN~\cite{yu2018stgcn} introduced spatio-temporal graph convolutional networks for traffic forecasting, while Graph WaveNet~\cite{wu2019graphwavenet} advanced deep spatial-temporal graph modeling. Notably, DCRNN-style encoders have also been adapted for ETA prediction in production settings.

Most recently, DuETA~\cite{dueta2023} introduced duration-aware ETA modeling at Baidu Maps, categorizing trips into short (0--3 km), medium (3--10 km), and long ($>$10 km) segments. DuETA achieved MAEs of 27s, 46s, and 98s respectively across these categories, significantly improving upon prior baselines. DuETA uses route information but focuses on aggregated road segments rather than individual vehicle trajectories with explicit route intent.

\subsection{Limitations of Existing Approaches}

However, existing benchmarks such as NYC~\cite{nyc_tlc}, Porto~\cite{moreira2013porto}, Chengdu/DiDi~\cite{didi2016}, and Geolife~\cite{zheng2012geolife} lack explicit representation of pre-planned routes. Models must therefore infer likely paths between origin and destination, introducing ambiguity and reducing accuracy. Most ETA prediction methods either use static sensor networks (e.g., DCRNN~\cite{dcrnn2018}, ST-GCN~\cite{yu2018stgcn}) that aggregate traffic at fixed locations, or rely on trajectory-based learning (e.g., DeepTTE~\cite{deepTTE2018}, STAD~\cite{abbar2020stad}) that infers routes implicitly from origin-destination pairs. Current state-of-the-art methods such as DuETA~\cite{dueta2023} use route information but focus on aggregated road segments rather than individual vehicle trajectories with explicit route intent.

Beyond methodological limitations, existing datasets present significant constraints that hinder comprehensive model comparison and evaluation. Public datasets such as NYC Taxi~\cite{nyc_tlc} and Porto Taxi~\cite{moreira2013porto} provide trip-level records with origin-destination pairs and GPS trajectories, but lack explicit route information and graph structures. These datasets require models to infer routes implicitly from origin-destination pairs, introducing ambiguity and reducing accuracy. Moreover, proprietary datasets used by production systems, such as the Baidu Maps dataset employed by DuETA~\cite{dueta2023}, are not publicly accessible, preventing independent validation and fair comparison across different model architectures.

The absence of a unified dataset that supports all model types creates a fundamental challenge for comprehensive evaluation. Trajectory-based models (DeepTTE~\cite{deepTTE2018}, STAD~\cite{abbar2020stad}) require GPS trajectory sequences, static graph models (DCRNN~\cite{dcrnn2018}, ST-GCN~\cite{yu2018stgcn}) require aggregated sensor network data, and route-aware models (DuETA~\cite{dueta2023}) require route segment information. No existing public dataset provides the combination of explicit routes, dynamic graph structure, and vehicle-level data required for evaluating all model types simultaneously. This limitation necessitates the creation of a simulated dataset with multiple variants, where each variant is tailored to a specific model's requirements while maintaining the same underlying traffic simulation, enabling fair comparison across all models.

These limitations highlight the need for approaches that: (1) model traffic at the vehicle level rather than aggregated sensor networks, (2) explicitly encode predefined routes rather than inferring them implicitly, (3) combine dynamic graph structures with temporal learning to capture evolving traffic conditions, and (4) provide a unified evaluation framework that enables fair comparison across different model architectures through a multi-variant dataset generated from the same underlying traffic simulation.


